\section{Assignment and Preservation}
\label{sec:newmethods}
\paragraph{A fixed candidate pool.}
Each training image has five source captions and hypotheses from e-SNLI-VE \citep{kayser2021evil}.
Hypotheses carry supported, contradicted, or neutral labels.
Every batch contains the same images, source captions, and hypotheses across policies.
Unpromoted hypotheses remain competitors.
Thus the intervention changes positive assignment without changing the available strings.

\emph{Source} uses only source captions as positives.
\emph{Supported}, denoted U, also promotes supported hypotheses for their annotated image.
The \emph{score-stratified} control promotes the same number within each of two frozen-similarity groups.
It draws assignments independently three times and keeps each assignment fixed during training.

\paragraph{Two effects of positive expansion.}
Let $s_{ij}$ score image $i$ against caption $j$.
Let $P_i$ contain its positive captions.
The image-query loss is
\begin{equation}
\ell_i=\log\sum_j e^{s_{ij}}-\frac{1}{|P_i|}\sum_{j\in P_i}s_{ij}.
\label{eq:newrow}
\end{equation}
With $K$ source captions and $M$ promoted hypotheses, each source target falls from $1/K$ to $1/(K+M)$.
The reverse loss averages caption queries against the image candidates.
Expansion adds eligible hypothesis queries and reduces the source queries' share of that average.
The symmetric loss gives each direction weight $1/2$.
Matching promotion counts preserves both changes across supported and randomized assignments.
Appendix~\ref{app:followupinvariants} proves the resulting source-column gradient identity at fixed scores.
We also cross Source and Supported supervision independently in the two loss directions.
The four cells measure each directional intervention and their interaction at matched schedules.

\paragraph{Allocation and distillation as separate interventions.}
Let $L_S$ and $L_U$ denote the complete Source and Supported losses.
Both losses use the unchanged candidate pool.
We compare the family
\begin{equation}
L_{\lambda,\beta}
=\lambda L_S+(1-\lambda)L_U+\beta T^2D_S.
\label{eq:adobjective}
\end{equation}
Here $\lambda$ controls source allocation and $\beta$ controls preservation.
$D_S$ is the mean teacher-to-student KL divergence on source-caption columns in both directions.
The frozen encoder supplies the teacher probabilities at temperature $T=2$.
Hypotheses receive supervised gradients but do not enter the distillation target.

The four cells separate the two interventions:
U uses $(\lambda,\beta)=(0,0)$;
A uses $(\lambda,0)$;
D uses $(0,\beta)$;
AD uses $(\lambda,\beta)$.
Allocation mixes complete objectives.
For example, the image-row source mass becomes $\lambda+(1-\lambda)K/(K+M)$.

\paragraph{Fair preservation controls.}
We also evaluate Source and WiSE-FT interpolation of Supported weights \citep{wortsman2022wise}.
All procedures receive the same development objective and epoch choices.
AD receives a score-stratified randomized control with identical coefficients, learning rate, and selected epochs.

The paired four-cell analysis uses the AD-selected schedule for U, A, D, and AD.
For any outcome $Y$, its interaction is
\begin{equation}
\mathcal I_Y=[Y(AD)-Y(D)]-[Y(A)-Y(U)].
\label{eq:interaction}
\end{equation}
An independently selected AD-versus-D comparison instead evaluates complete procedures.
